% ═══════════════════════════════════════════════════════════════
% MUSTERLÖSUNG - DS-03: Repaso general
% Klasse 9, Unidad 1+2 - Klausurvorbereitung
% ═══════════════════════════════════════════════════════════════

\documentclass[10pt, a4paper]{article}

\usepackage[utf8]{inputenc}
\usepackage[T1]{fontenc}
\usepackage[ngerman]{babel}

\usepackage{palatino}
\usepackage{helvet}
\newcommand{\sanstext}[1]{{\fontfamily{phv}\selectfont #1}}

\usepackage{microtype}
\usepackage[margin=1.4cm, top=1.6cm, bottom=1.3cm]{geometry}
\usepackage{enumitem}
\usepackage{tabularx}
\usepackage{booktabs}
\usepackage{array}
\usepackage{multicol}
\usepackage{tikz}
\usetikzlibrary{calc}

\usepackage{fancyhdr}
\usepackage{lastpage}
\usepackage{needspace}

\usepackage{xcolor}
\usepackage{tcolorbox}
\tcbuselibrary{skins}

% ═══════════════════════════════════════════════════════════════
% FARBEN
% ═══════════════════════════════════════════════════════════════

\definecolor{kopfzeile}{HTML}{1A1A1A}
\definecolor{akzent}{HTML}{C62828}
\definecolor{hellgrau}{HTML}{F2F2F2}
\definecolor{mittelgrau}{HTML}{777777}
\definecolor{dunkelgrau}{HTML}{444444}
\definecolor{loesung}{HTML}{C62828}

% ═══════════════════════════════════════════════════════════════
% VARIABLEN
% ═══════════════════════════════════════════════════════════════

\newcommand{\thema}{Repaso general -- Klausurvorbereitung}
\newcommand{\klasse}{9}
\newcommand{\maxpunkte}{30}

% ═══════════════════════════════════════════════════════════════
% KOPF- UND FUSSZEILE
% ═══════════════════════════════════════════════════════════════

\pagestyle{fancy}
\fancyhf{}
\renewcommand{\headrulewidth}{0pt}
\renewcommand{\footrulewidth}{0pt}
\cfoot{\sanstext{\footnotesize\textcolor{mittelgrau}{\thepage\,/\,\pageref{LastPage}}}}

% ═══════════════════════════════════════════════════════════════
% BOXEN
% ═══════════════════════════════════════════════════════════════

\newtcolorbox{stationbox}[1][]{
    colback=hellgrau,
    colframe=hellgrau,
    boxrule=0pt,
    arc=0pt,
    left=8pt, right=8pt, top=6pt, bottom=6pt,
    borderline north={2pt}{0pt}{akzent}
}

% ═══════════════════════════════════════════════════════════════
% BEFEHLE
% ═══════════════════════════════════════════════════════════════

\newcommand{\es}[1]{\textbf{#1}}
\newcommand{\de}[1]{{\small\textit{\textcolor{mittelgrau}{#1}}}}
\newcommand{\luecke}[1]{\rule{#1}{0.4pt}}
\newcommand{\punktebox}[1]{\hfill\sanstext{\small[\_\_/#1]}}

% Lösung in Rot
\newcommand{\lsg}[1]{\textcolor{loesung}{\textbf{#1}}}

\newcommand{\aufgabe}[2]{%
    \needspace{4\baselineskip}%
    \vspace{10pt}%
    \noindent\sanstext{\textbf{#1}\quad#2}%
    \vspace{5pt}%
    \nopagebreak[4]%
}

\newlist{teil}{enumerate}{1}
\setlist[teil]{label=\alph*), leftmargin=1.8em, itemsep=3pt, parsep=0pt, topsep=3pt}

\setlength{\parindent}{0pt}
\setlength{\parskip}{4pt}
\setlength{\columnsep}{20pt}

% ═══════════════════════════════════════════════════════════════
% DOKUMENT
% ═══════════════════════════════════════════════════════════════

\begin{document}

% ═══════════════════════════════════════════════════════════════
% HEADER
% ═══════════════════════════════════════════════════════════════

\noindent
\begin{tikzpicture}[remember picture, overlay]
    \fill[kopfzeile] (current page.north west) rectangle +(\paperwidth, -1cm);
    \node[anchor=west, text=white, font=\sanstext\large\bfseries]
        at ([xshift=1.4cm, yshift=-0.5cm]current page.north west) {\thema{} -- \textcolor{loesung}{MUSTERLÖSUNG}};
    \node[anchor=east, text=white, font=\sanstext\small]
        at ([xshift=-1.4cm, yshift=-0.5cm]current page.north east)
        {Spanisch\,·\,Klasse \klasse\,·\,ML};
\end{tikzpicture}

\vspace{0.5cm}

\noindent
\sanstext{\small\textbf{Name:}} \luecke{5cm} \hfill
\sanstext{\small\textbf{Datum:}} \luecke{2.5cm}

\vspace{6pt}

% ═══════════════════════════════════════════════════════════════
% STATION 1: VOKABELN
% ═══════════════════════════════════════════════════════════════

\begin{stationbox}
\sanstext{\footnotesize\textbf{Station 1: Vocabulario U1 + U2}}
\end{stationbox}

\vspace{2pt}

\aufgabe{Aufgabe 1}{Vokabeltraining: Übersetze und ergänze.} \punktebox{6}

\vspace{2pt}

{\small\textbf{A) Übersetze ins Spanische:}}

\vspace{4pt}

\begin{multicols}{2}
{\footnotesize
\begin{enumerate}[leftmargin=1.5em, itemsep=4pt]
    \item der Kühlschrank: \lsg{la nevera}
    \item das Schlafzimmer: \lsg{el dormitorio}
    \item die Dusche: \lsg{la ducha}
    \item der Sessel: \lsg{el sillón}
    \item der Spiegel: \lsg{el espejo}
    \item das Regal: \lsg{la estantería}
\end{enumerate}
}
\end{multicols}

\vspace{4pt}

{\small\textbf{B) Ergänze die passenden Wörter:}}

\vspace{4pt}

{\footnotesize
\begin{enumerate}[leftmargin=1.5em, itemsep=4pt, start=7]
    \item En la \lsg{cocina} hay una nevera y un horno.
    \item El sofá está en el \lsg{salón}.
    \item Mi cama está en el \lsg{dormitorio}.
    \item El coche está en el \lsg{garaje} (die Garage).
    \item Voy al \lsg{baño} para ducharme.
    \item La lámpara está al lado de la \lsg{cama}.
\end{enumerate}
}

\vspace{8pt}

% ═══════════════════════════════════════════════════════════════
% STATION 2: PRETÉRITO PERFECTO
% ═══════════════════════════════════════════════════════════════

\begin{stationbox}
\sanstext{\footnotesize\textbf{Station 2: El pretérito perfecto}}
\end{stationbox}

\vspace{2pt}

\aufgabe{Aufgabe 2}{Konjugiere und bilde Sätze im pretérito perfecto.} \punktebox{8}

\vspace{2pt}

{\small\textbf{A) Ergänze die Konjugation von \es{haber}:}}

\vspace{4pt}

{\footnotesize
\begin{tabular}{@{}llllll@{}}
yo \lsg{he} & tú \lsg{has} & él/ella \lsg{ha} & nosotros \lsg{hemos} & vosotros \lsg{habéis} & ellos \lsg{han}
\end{tabular}
}

\vspace{6pt}

{\small\textbf{B) Schreibe das Partizip:}}

\vspace{4pt}

{\footnotesize
\begin{tabular}{@{}llll@{}}
hablar $\rightarrow$ \lsg{hablado} & comer $\rightarrow$ \lsg{comido} & vivir $\rightarrow$ \lsg{vivido} & poner $\rightarrow$ \lsg{puesto}
\end{tabular}
}

\vspace{6pt}

{\small\textbf{C) Bilde Sätze im pretérito perfecto:}}

\vspace{4pt}

\begin{teil}
    \item yo / hablar / con mi madre $\rightarrow$ \lsg{Yo he hablado con mi madre.}
    \item nosotros / comer / en un restaurante $\rightarrow$ \lsg{Nosotros hemos comido en un restaurante.}
    \item ellos / vivir / en Madrid $\rightarrow$ \lsg{Ellos han vivido en Madrid.}
    \item tú / hacer / los deberes $\rightarrow$ \lsg{Tú has hecho los deberes.}
\end{teil}

\vspace{8pt}

% ═══════════════════════════════════════════════════════════════
% STATION 3: BEWEGUNGSVERBEN + PRONOMEN
% ═══════════════════════════════════════════════════════════════

\begin{stationbox}
\sanstext{\footnotesize\textbf{Station 3: Verbos de movimiento + Pronombres}}
\end{stationbox}

\vspace{2pt}

\aufgabe{Aufgabe 3}{Ersetze das Nomen durch ein Pronomen und bilde Sätze.} \punktebox{6}

\vspace{2pt}

{\small Schreibe jeden Satz mit dem passenden Objektpronomen (\es{lo}, \es{la}, \es{los}, \es{las}) um.}

\vspace{4pt}

\begin{teil}
    \item \es{Pongo la lámpara en la mesa.} $\rightarrow$ \lsg{La pongo en la mesa.}
    \item \es{Subimos los libros al piso.} $\rightarrow$ \lsg{Los subimos al piso.}
    \item \es{Sacamos la nevera del camión.} $\rightarrow$ \lsg{La sacamos del camión.}
    \item \es{Llevo el televisor al salón.} $\rightarrow$ \lsg{Lo llevo al salón.}
    \item \es{Bajamos las cajas del coche.} $\rightarrow$ \lsg{Las bajamos del coche.}
    \item \es{Traigo la guitarra de mi habitación.} $\rightarrow$ \lsg{La traigo de mi habitación.}
\end{teil}

\vspace{8pt}

% ═══════════════════════════════════════════════════════════════
% STATION 4: LESEVERSTEHEN
% ═══════════════════════════════════════════════════════════════

\begin{stationbox}
\sanstext{\footnotesize\textbf{Station 4: Comprensión lectora}}
\end{stationbox}

\vspace{2pt}

\aufgabe{Aufgabe 4}{Lies den Text und beantworte die Fragen.} \punktebox{5}

\vspace{4pt}

{\small Beantworte auf Spanisch:}

\vspace{4pt}

\begin{teil}
    \item ¿Cuántas habitaciones tiene el piso?\\
    $\rightarrow$ \lsg{El piso tiene cuatro habitaciones.}

    \item ¿Por qué le gusta su habitación al narrador?\\
    $\rightarrow$ \lsg{Le gusta porque tiene un balcón.}

    \item ¿Qué ha hecho el padre?\\
    $\rightarrow$ \lsg{El padre ha puesto el sofá en el salón.}

    \item ¿Dónde ha puesto la madre la ropa?\\
    $\rightarrow$ \lsg{La madre ha puesto la ropa en el armario.}

    \item ¿Qué van a comprar esta tarde?\\
    $\rightarrow$ \lsg{Van a comprar una lámpara nueva para el salón.}
\end{teil}

\vspace{8pt}

% ═══════════════════════════════════════════════════════════════
% STATION 5: TEXTPRODUKTION
% ═══════════════════════════════════════════════════════════════

\begin{stationbox}
\sanstext{\footnotesize\textbf{Station 5: Producción de textos}}
\end{stationbox}

\vspace{2pt}

\aufgabe{Aufgabe 5}{Schreibe einen kurzen Text über dein Zimmer.} \punktebox{5}

\vspace{2pt}

{\footnotesize\textbf{Mögliche Lösung:}}

\vspace{4pt}

{\small\textcolor{loesung}{
Mi habitación es pequeña pero muy bonita. Está en el primer piso de nuestra casa. En mi habitación hay una cama grande, un armario y una estantería con muchos libros. También tengo un escritorio donde hago los deberes. Me gusta mi habitación porque tiene una ventana grande y hay mucha luz. Esta semana he puesto un póster nuevo en la pared.
}}

% ═══════════════════════════════════════════════════════════════
% FOOTER
% ═══════════════════════════════════════════════════════════════

\vfill

\noindent\rule{\textwidth}{0.5pt}

\vspace{4pt}

\hfill\sanstext{\textbf{Punkte:}} \luecke{1.2cm} / \maxpunkte

\end{document}
